\subsection{谱与预解式}

引理 4. --- 设 E 为复巴拿赫空间，u 为 E 上的闭单射部分算子，满足 \(u^{-1} \in \mathcal{L}(\mathrm{E})\)。设 \(v \in \mathcal{L}(\mathrm{E})\)

满足 \(\|v\| < \|u^{-1}\|^{-1}\)。则部分算子 \(u + v\) 是单射的，我们有 \((u + v)^{-1} \in \mathcal{L}(\mathrm{E})\)，并且

\[
(u + v) ^ {- 1} = u ^ {- 1} \circ \sum_ {k = 0} ^ {+ \infty} (- v u ^ {- 1}) ^ {k},\tag{6}
\]

其中该级数在 \(\mathcal{L}(\mathrm{E})\) 中绝对收敛。此外，我们有

\[
\| (u + v) ^ {- 1} \| \leqslant \frac {\| u ^ {- 1} \|}{1 - \| v \| \| u ^ {- 1} \|}.
\]

由于 \(\|v\|\|u^{-1}\|<1\)，通项为 \((-vu^{-1})^{k}\) 的级数在 \(\mathcal{L}(\mathrm{E})\) 中绝对收敛，且其和是自同态 \(1_{E}+vu^{-1}\) 的逆（I，第 22 页，命题 2）。于是，部分算子 \((1_{\mathrm{E}}+vu^{-1})\circ u=u+v\)（IV，第 226 页，注 vi）是单射的，且逆部分算子 \((u+v)^{-1}=u^{-1}\circ(1_{\mathrm{E}}+vu^{-1})^{-1}\)（IV，第 226 页，注 iii）属于 \(\mathcal{L}(\mathrm{E})\)。由于

\[
\left\| \left(1 _ {\mathrm{E}} + v u ^ {- 1}\right) ^ {- 1} \right\| \leqslant \sum_ {k = 0} ^ {+ \infty} (\left\| v \right\| \left\| u ^ {- 1} \right\|) ^ {k} = \frac {1}{1 - \left\| v \right\| \left\| u ^ {- 1} \right\|},
\]

引理得证。

定义 8. --- 设 u 是复巴拿赫空间 E 上的稠定闭算子。u 的预解集是由这样的复数 λ 组成的集合：部分算子 \(\lambda 1_{\mathrm{E}} - u\) 是单射的，且其逆 \((\lambda 1_{\mathrm{E}} - u)^{-1}\) 属于 \(\mathcal{L}(\mathrm{E})\)。

u 的谱，记作 \(\mathrm{Sp}(u)\)，是预解集在 C 中的余集。

若 \(\lambda \in \mathbf{C} - \mathrm{Sp}(u)\)，我们用 \(\mathrm{R}(u,\lambda) \in \mathcal{L}(\mathrm{E})\) 表示 \(\lambda 1_{\mathrm{E}} - u\) 的逆。从 \(\mathbf{C} - \mathrm{Sp}(u)\) 到 \(\mathcal{L}(\mathrm{E})\) 中把 \(\lambda\) 映为 \(\mathrm{R}(u,\lambda)\) 的映射称为 \(u\) 的预解式。

注记。--- 设 E 为复巴拿赫空间，u 为 E 上稠定的闭算子。

\begin{enumerate}
\def\labelenumi{\arabic{enumi})}
\item
  若 \(u \in \mathcal{L}(\mathrm{E})\)，则它的谱与元素 \(u\) 作为代数 \(\mathcal{L}(\mathrm{E})\) 中元素的谱一致（I，第 2 页，定义 1）。
\item
  设 D 为 \(u\) 的定义域。对每个 \(\lambda \in \mathbf{C} - \mathrm{Sp}(u)\)，我们有
\end{enumerate}

\[
1 _ {\mathrm{E}} = (\lambda 1 _ {\mathrm{E}} - u) \circ \mathrm{R} (u, \lambda), \quad 1 _ {\mathrm{D}} = \mathrm{R} (u, \lambda) \circ (\lambda 1 _ {\mathrm{E}} - u).\tag{7}
\]

此外，对于 u 的预解集中的 \(\lambda_{1}\) 与 \(\lambda_{2}\)，我们有

\[
\mathrm{R} (u, \lambda_ {1}) - \mathrm{R} (u, \lambda_ {2}) = (\lambda_ {2} - \lambda_ {1}) \mathrm{R} (u, \lambda_ {2}) \circ \mathrm{R} (u, \lambda_ {1}),\tag{8}
\]

\[
\mathrm{R} (u, \lambda_ {1}) \circ \mathrm{R} (u, \lambda_ {2}) = \mathrm{R} (u, \lambda_ {2}) \circ \mathrm{R} (u, \lambda_ {1}),\tag{9}
\]

事实上，有

\[
\mathrm{R} (u, \lambda_ {1}) - \mathrm{R} (u, \lambda_ {2}) = \mathrm{R} (u, \lambda_ {1}) \circ 1 _ {\mathrm{E}} - 1 _ {\mathrm{D}} \circ \mathrm{R} (u, \lambda_ {2}).
\]

由于 \(1_{\mathrm{E}} = (\lambda_21_{\mathrm{E}} - u)\circ \mathrm{R}(u,\lambda_2)\) 与 \(1_{\mathrm{D}} = \mathrm{R}(u,\lambda_1)\circ (\lambda_11_{\mathrm{E}} - u)\)，我们得到

\[
\begin{array}{c} \mathrm{R} (u, \lambda_ {1}) - \mathrm{R} (u, \lambda_ {2}) = \lambda_ {2} \mathrm{R} (u, \lambda_ {1}) \circ \mathrm{R} (u, \lambda_ {2}) \\ - \mathrm{R} (u, \lambda_ {1}) \circ u \circ \mathrm{R} (u, \lambda_ {2}) - \lambda_ {1} \mathrm{R} (u, \lambda_ {1}) \circ \mathrm{R} (u, \lambda_ {2}) \\ + \mathrm{R} (u, \lambda_ {1}) \circ u \circ \mathrm{R} (u, \lambda_ {2}), \end{array}
\]

由此得到第一个公式。交换 \(\lambda_{1}\) 与 \(\lambda_{2}\) 的角色，便推出第二个公式。

\begin{enumerate}
\def\labelenumi{\arabic{enumi})}
\setcounter{enumi}{2}
\item
  设 \(\lambda \in \mathbf{C}\)。由闭图像定理（EVT，I，第 19 页，推论 5），若从 \(\mathrm{dom}(u)\) 到 E 中的由 \(x \mapsto (\lambda 1_{\mathrm{E}} - u)(x)\) 定义的线性映射是双射，则它的逆（其图像是闭的）从 E 到 E 连续。因此，\(\lambda\) 属于 \(u\) 的预解集，当且仅当由 \(x \mapsto (\lambda 1_{\mathrm{E}} - u)(x)\) 给出的从 \(\mathrm{dom}(u)\) 到 E 中的映射是双射。
\item
  若 \(\lambda\) 属于 u 的谱，则下列性质之一成立：
\end{enumerate}

\begin{enumerate}
\def\labelenumi{(\roman{enumi})}
\item
  \(\lambda 1_{\mathrm{E}} - u\) 的核不化为 0；这时称 \(\lambda\) 是 \(u\) 的特征值，并称 \(\operatorname{Ker}(\lambda 1_{\mathrm{E}} - u)\) 的维数为其重数；
\item
  部分算子 \(\lambda 1_{\mathrm{E}} - u\) 是单射的，而其像在 E 中不稠密；这时称 \(\lambda\) 属于 \(u\) 的剩余谱；
\item
  部分算子 \(\lambda 1_{\mathrm{E}} - u\) 是单射的，其像在 E 中稠密，但 \(\lambda 1_{\mathrm{E}} - u\) 不是满射；这时称 \(\lambda\) 属于 \(u\) 的连续谱。
\end{enumerate}

\begin{enumerate}
\def\labelenumi{\arabic{enumi})}
\setcounter{enumi}{4}
\tightlist
\item
  设 \(\lambda\) 是属于 \(u\) 的预解集的复数。预解式 \(\mathrm{R}(u,\lambda)\) 是从 E 到 E 中的单射线性映射；其像是 \(u\) 的定义域，且 \(u = \lambda 1_{\mathrm{E}} - \mathrm{R}(u,\lambda)^{-1}\)（参见 IV，第 226 页）。反之，这一性质刻画了预解集与预解式。准确地说，设 \(\lambda \in \mathbf{C}\)；若存在从 E 到 E 中的连续单射线性映射 w，使得 \(u = \lambda 1_{E} - w^{-1}\)，则 \(\lambda\) 属于 u 的预解集，且 \(w = \mathrm{R}(u, \lambda)\)。
\end{enumerate}

特别地，若 v 是 E 上的闭稠定算子，且 \(\lambda \in C\) 是不属于 \(\operatorname{Sp}(u) \cup \operatorname{Sp}(v)\) 的复数，则等式 \(\operatorname{R}(u, \lambda) = \operatorname{R}(v, \lambda)\) 蕴含 u = v。

\begin{enumerate}
\def\labelenumi{\arabic{enumi})}
\setcounter{enumi}{5}
\tightlist
\item
  我们把可闭部分算子 \(u\) 的谱定义为其闭包的谱。
\end{enumerate}

存在谱为空集、或谱等于 C 的闭算子（IV，第 347 页，习题 12）。

设 E 为复巴拿赫空间，u 为 E 上的闭稠定算子。若 F 是复巴拿赫空间，且 \(v: E \to F\) 是同构，则我们有 \(\mathrm{Sp}(v \circ u \circ v^{-1}) = \mathrm{Sp}(u)\) 与 \(\mathrm{R}(v \circ u \circ v^{-1}, \lambda) = v \circ \mathrm{R}(u, \lambda) \circ v^{-1}\)，对一切 \(\lambda \notin \mathrm{Sp}(u)\) 成立。

命题 14. --- 设 E 为复巴拿赫空间。设 u 为 E 上稠定的闭算子，U 为其预解集。a) 对每个 \(\lambda\in\mathrm{U}\)，C 中以 \(\lambda\) 为心、以 \(\|\mathrm{R}(u,\lambda)\|^{-1}\) 为半径的开圆盘都包含于 U；

\begin{enumerate}
\def\labelenumi{\alph{enumi})}
\setcounter{enumi}{1}
\item
  集合 U 是 C 中的开集；
\item
  设 \(\operatorname{Sp}(u)\) 非空。设 \(\lambda \in \mathrm{U}\)，\(\delta\) 是 \(\mathbf{C}\) 中从 \(\lambda\) 到 \(u\) 的谱的距离。我们有 \(\delta > 0\) 及 \(\| \mathrm{R}(u, \lambda) \| \geqslant 1 / \delta\)；
\item
  映射 \(\lambda \mapsto \mathrm{R}(u,\lambda)\) 是从 U 到 \(\mathcal{L}(\mathrm{E})\) 中的全纯映射。对每个整数 \(k\in \mathbf{N}\) 与每个 \(\lambda \in \mathrm{U}\)，我们有
\end{enumerate}

\[
\frac {\partial^ {k}}{\partial \lambda^ {k}} \mathrm{R} (u, \lambda) = (- 1) ^ {k} k! \mathrm{R} (u, \lambda) ^ {k + 1}.
\]

设 \(\lambda \in \mathrm{U}\)。对每个 \(\mu \in \mathbf{C}\)，只要 \(\| (\mu -\lambda)1_{\mathrm{E}}\| < \| \mathrm{R}(u,\lambda)\|^{-1}\)，将引理 4 应用于单射部分算子 \(\lambda 1_{\mathrm{E}} - u\) 并取 \(v = (\mu -\lambda)1_{\mathrm{E}}\)，便知部分算子 \(\mu 1_{\mathrm{E}} - u = \lambda 1_{\mathrm{E}} - u + v\) 是单射的且具有连续的逆。这就蕴含 \(a\)）。由同上所引，我们还有公式

\[
\mathrm{R} (u, \mu) = \mathrm{R} (u, \lambda) \circ \sum_ {k \in \mathbf {N}} (\lambda - \mu) ^ {k} \mathrm{R} (u, \lambda) ^ {k},
\]

由此推出 u 的预解式在 U 中全纯。

断言 \(b\)）由 \(a\)）立得。若 \(\mathrm{Sp}(u)\) 非空，则 \(\lambda\) 到 \(\mathrm{Sp}(u)\) 的距离严格为正（TG，IX，第 13 页，命题 2），由此得 \(c\)）。

断言 \(d\)）的最后一部分对 \(k\) 用归纳法证明，其中 \(k = 1\) 的情形是第 245 页公式 (8) 的推论。

命题 15. --- 设 u 是复巴拿赫空间 E 上稠定的闭算子，λ 是属于 u 的预解集的复数。

\begin{enumerate}
\def\labelenumi{\alph{enumi})}
\item
  子集 \(\mathrm{Sp}(\mathrm{R}(u,\lambda)) \setminus \{0\}\)（\(\mathbf{C}\) 中的）是 \(u\) 的谱经映射 \(\mu \mapsto (\lambda -\mu)^{-1}\)（从 \(\mathbf{C} \setminus \{\lambda \}\) 到 \(\mathbf{C}^*\) 中）所成的像；
\item
  对每个 \(\mu \neq \lambda\)（在 \(\mathbf{C}\) 中），有
\end{enumerate}

\[
\operatorname{Ker} \left(\mu 1 _ {\mathrm{E}} - u\right) = \operatorname{Ker} \left(\left(\lambda - \mu\right) ^ {- 1} 1 _ {\mathrm{E}} - \mathrm{R} (u, \lambda)\right).
\]

我们来证明断言 \(a\)）。对每个 \(\mu \neq \lambda\)，计算

\[
\mu 1 _ {\mathrm{E}} - u = (\lambda - \mu) ((\lambda - \mu) ^ {- 1} 1 _ {\mathrm{E}} - \mathrm{R} (u, \lambda)) (\lambda 1 _ {\mathrm{E}} - u).
\]

由于 \(\lambda \notin \mathrm{Sp}(u)\) 且 \(\mu \neq \lambda\)，线性映射 \((\lambda - \mu)(\lambda 1_{\mathrm{E}} - u)\) 是从 \(\mathrm{dom}(u)\) 到 E 上的双射。于是，该公式表明：\(\mu 1_{\mathrm{E}} - u\) 是从 \(\mathrm{dom}(u)\) 到 E 上的双射，当且仅当 \((\lambda - \mu)^{-1} 1_{\mathrm{E}} - \mathrm{R}(u, \lambda)\) 是从 E 到 E 上的双射，由此即得断言。

我们来证明 \(b\)）。若 \(\mu \neq \lambda\) 且 \(x \in \mathrm{Ker}((\lambda - \mu)^{-1}1_{\mathrm{E}} - \mathrm{R}(u, \lambda))\)，则 \(x \in \mathrm{dom}(u)\)，且公式 \(1_{\mathrm{E}} = (\lambda 1_{\mathrm{E}} - u) \circ \mathrm{R}(u, \lambda)\) 蕴含 \(x \in \mathrm{Ker}(\mu 1_{\mathrm{E}} - u)\)。反之，若 \(x \in \mathrm{Ker}(\mu 1_{\mathrm{E}} - u)\) 且 \(\mu \neq \lambda\)，则公式 \(1_{\mathrm{dom}(u)} = \mathrm{R}(u, \lambda) \circ (\lambda 1_{\mathrm{E}} - u)\) 蕴含 \(\mathrm{R}(u, \lambda)(x) = (\lambda - \mu)^{-1}x\)。

命题 16. --- 设 u 是复希尔伯特空间 E 上稠定的闭算子。\(u^{*}\) 的谱是 u 的谱在复共轭下的像，且对 u 的预解集的每个元素 \(\lambda\)，有 \(\mathrm{R}(u,\lambda)^{*}=\mathrm{R}(u^{*},\overline{\lambda})\)。特别地，若 u 自伴，则自同态 \(\mathrm{R}(u,\lambda)\) 对每个 \(\lambda\notin\mathrm{Sp}(u)\) 都是正规的。

设 \(\lambda\in\mathbf{C} - \mathrm{Sp}(u)\) 是 \(u\) 的预解集中的一个元素。我们有 \(u=\lambda1_{\mathrm{E}}-\mathrm{R}(u,\lambda)^{-1}\)，于是

\[
u ^ {*} = \overline {{\lambda}} 1 _ {\mathrm{E}} - (\mathrm{R} (u, \lambda) ^ {- 1}) ^ {*} = \overline {{\lambda}} 1 _ {\mathrm{E}} - (\mathrm{R} (u, \lambda) ^ {*}) ^ {- 1}
\]

（IV，第 236 页及 IV，第 239 页命题 9）。由注 5 推出 \(\overline{\lambda} \in \mathbf{C} - \mathrm{Sp}(u^{*})\) 且 \(\mathrm{R}(u, \lambda)^{*} = \mathrm{R}(u^{*}, \overline{\lambda})\)。因此，\(u^{*}\) 的谱包含于 \(\mathrm{Sp}(u)\) 经复共轭所成的像。把该性质应用于 \(u^{*}\) 即得等式，因为 \(u^{**} = u\)（IV，第 237 页命题 8 的推论）。最后的断言则由第 245 页公式 (9) 得出。

推论。--- 设 u 是复希尔伯特空间 E 上的自伴部分算子。若 E 非零，则 u 的谱非空。

设 \(\operatorname{Sp}(u)\) 为空集。则 u 是单射的，且 \(u^{-1} = -\operatorname{R}(u,0)\) 是 E 的单射自同态，满足 \(\operatorname{Sp}(u^{-1}) \subset \{0\}\)（命题 15，a）），从而 \(\operatorname{Sp}(u^{-1}) = \{0\}\)（I，第 26 页，推论 1）。由于 \(u^{-1}\) 正规（命题 16），这就蕴含 \(u^{-1}\) 为零（I，第 110 页，例 1），此为矛盾。

引理 5. --- 设 E 为复希尔伯特空间，u 为 E 上的闭稠定部分算子。设 \(\lambda \in \mathbf{C}\)。假设存在实数 \(c > 0\)，使得

\[
\| u (x) - \lambda x \| \geqslant c \| x \| \quad \text{对一切 } x \in \operatorname{dom} (u),\tag{10}
\]

\[
\| u ^ {*} (x) - \overline {{\lambda}} x \| \geqslant c \| x \| \quad \text{对一切 } x \in \operatorname{dom} (u ^ {*}),\tag{11}
\]

则 \(\lambda\) 属于 \(u\) 的预解集，且 \(\| \mathrm{R}(u,\lambda)\| \leqslant c^{-1}\)。

假设蕴含 \(u - \lambda 1_{\mathrm{E}}\) 与 \(u^{*} - \overline{\lambda} 1_{\mathrm{E}}\) 都是单射的。设 F 为 \(u - \lambda 1_{\mathrm{E}}\) 的像。空间 F 在 E 中稠密，因为其正交补等于 \(\operatorname{Ker}(u^{*} - \overline{\lambda} 1_{\mathrm{E}})\)（IV，第 236 页，命题 7，c）），而后者为零。

我们来证明空间 F 是闭的。设 \((x_{n})_{n\in\mathbf{N}}\) 是 dom(u) 中的序列，使序列 \((u(x_{n})-\lambda x_{n})_{n\in\mathbf{N}}\) 收敛于 \(y\in F\)。不等式 (10) 蕴含序列 \((x_{n})_{n\in\mathbf{N}}\) 是 E 中的柯西序列。设 \(x\in E\) 为其极限。序列 \((x_{n},u(x_{n}))\) 在 E×E 中收敛于 \((x,y+\lambda x)\)；由于 u 的图像是闭的，我们因此有 \(x\in\text{dom}(u)\) 与 \(u(x)=y+\lambda x\)，这就证明了 \(y\in F\)。

由此断定 F = E。于是部分算子 \(u - \lambda1_{E}\) 是双射，从而 \(\lambda \notin \operatorname{Sp}(u)\)（注 3）。不等式 (10) 进而又蕴含 \(\|R(u, \lambda)\| \leqslant c^{-1}\)。

命题 17. — 设 E 为复希尔伯特空间，u 为 E 上的自伴部分算子。

\begin{enumerate}
\def\labelenumi{\alph{enumi})}
\item
  u 的谱包含于 \(\mathbf{R}\)；
\item
  若 u 是正的，则 u 的谱包含于 \(\mathbf{R}_{+}\)；
\item
  设 E 非零。设 \(\lambda \notin \operatorname{Sp}(u)\)，且 \(\delta > 0\) 为 \(\lambda\) 到 \(u\) 的谱的距离。我们有 \(\| \mathrm{R}(u, \lambda) \| = \delta^{-1}\)。
\end{enumerate}

设 \((a,b)\in\mathbf{R}\times\mathbf{R}\)，\(\lambda=a+ib\)。设 \(x\in\mathrm{dom}(u)\)。由于 u 自伴，我们有 \(\langle x|u(x)\rangle\in\mathbf{R}\)，由此

\[
\begin{array}{r} \| u (x) - \lambda x \| ^ {2} = \| u (x) \| ^ {2} - 2 a \langle x | u (x) \rangle + (a ^ {2} + b ^ {2}) \| x \| ^ {2} \\ = \| u (x) - \overline {{\lambda}} x \| ^ {2}. \end{array}
\]

设 \(b \neq 0\)。于是得到

\[
\begin{array}{c} \| u (x) - \lambda x \| ^ {2} = \| u (x) - \overline {{\lambda}} x \| ^ {2} \geqslant (\| u (x) \| - a \| x \|) ^ {2} + b ^ {2} \| x \| ^ {2} \\ \geqslant b ^ {2} \| x \| ^ {2}. \end{array}
\]

于是由引理 5 得 \(\lambda \notin \operatorname{Sp}(u)\)，即断言 \(a\)）。

设 u 还是正的。若 b = 0 且 a \textless{} 0，则对 \(x \in \text{dom}(u)\) 类似地得到不等式

\[
\left\| u (x) - \lambda x \right\| ^ {2} = \left\| u (x) - \overline {{\lambda}} x \right\| ^ {2} \geqslant \left(\left\| u (x) \right\| - a \| x \|\right) ^ {2} \geqslant a ^ {2} \| x \| ^ {2},
\]

从而 \(\lambda\notin\mathrm{Sp}(u)\)（同上所引），这就证明了 \(b\)）。

最后证明 \(c\)）。由命题 16，预解式 \(\mathrm{R}(u,\lambda)\) 是 E 的正规自同态。因此其范数等于其谱半径（I，第 108 页，推论 1），由此

\[
\| \mathrm{R} (u, \lambda) \| = \sup _ {\mu \in \operatorname{Sp} (\mathrm{R} (u, \lambda))} | \mu |
\]

（I，第 24 页，定理 1）。R(u,λ) 的谱不可能化为 \{0\}，因为否则将有 \(\| \mathrm{R}(u,\lambda)\| = 0\)，于是作为 R(u,λ) 之像的 \(\operatorname{dom}(u)\) 将为零，E 亦然。因此命题 15 蕴含

\[
\| \mathrm{R} (u, \lambda) \| = \sup _ {\mu \in \mathrm{Sp} (u)} \frac {1}{| \lambda - \mu |} = \frac {1}{\delta}.
\]

推论。--- 设 u 是 E 上的自伴算子。

\begin{enumerate}
\def\labelenumi{\alph{enumi})}
\item
  u 的剩余谱是空集；
\item
  对所有 \(\lambda \neq \mu\)（在 \(\mathbf{C}\) 中），\(u\) 的属于 \(\lambda\) 与 \(\mu\) 的特征空间相互正交。
\end{enumerate}

我们来证明 a)。设 \(\lambda\) 属于 \(u\) 的谱；它是实数（命题 17）。我们有 \(\mathrm{Ker}(\lambda 1_{\mathrm{E}} - u) = \mathrm{Im}(\lambda 1_{\mathrm{E}} - u)^{\circ}\)（IV，第 236 页，命题 7），因此 \(\lambda\) 不是 \(u\) 的特征值，只要部分算子 \(\lambda 1_{\mathrm{E}} - u\) 的像稠密。按定义，这就蕴含 \(u\) 的剩余谱是空集。

我们来证明 \(b\)）。由该命题，可以假设 \(\lambda\) 与 \(\mu\) 是实数。设 \(x \in \operatorname{dom}(u)\) 满足 \(u(x) = \lambda x\)，且 \(y \in \operatorname{dom}(u)\) 满足 \(u(y) = \mu y\)。于是有

\[
\lambda \langle x \mid y \rangle = \langle u (x) \mid y \rangle = \langle x \mid u (y) \rangle = \mu \langle x \mid y \rangle ,
\]

由此 \(\langle x \mid y \rangle = 0.\)

注。--- 若 u 是非自伴的闭对称算子，则其谱不包含于 R（参见下文 IV，第 257 页推论 10），而且 u 的对应于 λ 与 \(\overline{\lambda}\) 的特征空间有可能不正交（IV，第 347 页，习题 11）。

